% ============================================================================
%  Documento de costos frente a venta. Informe 3.
%  Formulario T-22 (FEP01 p.69): proveedores clave, adquisiciones clave,
%  curva S, análisis de costos, VAN y TIR, y valorización de las cinco
%  innovaciones en el flujo de caja. Las cifras salen de
%  economico/datos/partidas.csv: aquí no se escribe ningún monto a mano.
% ============================================================================

\begin{resumenApertura}
El documento compara los costos de implementación y operación con el precio de la solución para Transportes Curimón. Presenta las adquisiciones, el calendario financiero y los indicadores del proyecto.
\begin{recibe}
  \item Desglose de costos y precios por fase, respaldado en la planilla de cálculo.
\end{recibe}
\end{resumenApertura}

\section{Proveedores clave de la solución}\label{sec:eco-proveedores}

La \tabref{eco-proveedores} identifica proveedores relevantes para el despliegue, la continuidad del servicio y la adquisición del equipamiento de referencia.

\begin{tabla}{Proveedores clave de la solución}{eco-proveedores}{L{38mm} X L{38mm} L{34mm}}
Proveedor & Qué provee & Componente o servicio & Criterio de selección \\ \midrule
Microsoft & Azure en Chile Central y Brazil South & Plataforma y recuperación & Disponibilidad regional y canal CSP \\
Dell & Servidores & Infraestructura on-premise & Soporte y disponibilidad en Chile \\
iWave, Teltonika y Technoton & Equipos y periféricos de campo y red & Telemetría y acceso en terminal & Compatibilidad con el diseño de T-11 \\
Fortinet & Equipos de seguridad y red & Sala técnica y terminales & Compatibilidad con la arquitectura y soporte local \\
GitLab y SonarSource & Licencias de desarrollo y análisis de código & Integración continua y calidad & Cobertura de los procesos de desarrollo y pruebas \\
Ground Control & Tiempo de aire Iridium & Respaldo satelital & Cobertura en rutas sin servicio terrestre \\
\end{tabla}

El operador móvil y el proveedor del circuito ExpressRoute se seleccionan al validar cobertura y factibilidad del enlace.

\section{Adquisiciones clave o relevantes}\label{sec:eco-adquisiciones}

La \tabref{eco-adquisiciones} presenta las compras y servicios que condicionan la habilitación, el piloto y la continuidad. Los meses corresponden al calendario presupuestado. Las fechas de orden y entrega se cierran con cada proveedor.

\begin{tabla}[ancho=completo]{Adquisiciones clave y dependencias}{eco-adquisiciones}{L{33mm} L{17mm} L{19mm} X X}
Adquisición & Pagador & Mes previsto & Fecha requerida & Plazo, condición o alternativa \\ \midrule
Grupo electrógeno & audIT & M1 & H3, M6 & Hasta seis meses de plazo. Confirmar entrega compatible con H3. \\
Sala técnica, clima y extinción & audIT & M3 & H3, M6 & Cotizar instalación certificada, disponibilidad y fecha de puesta en servicio. \\
Cómputo de terminal y red & audIT y Curimón & M3--M4 & H3, M6 & audIT compra cómputo y gabinetes. Curimón adquiere equipos de red. Plazo de referencia AAEON de una a tres semanas. \\
Azure y conectividad privada & audIT & M1, M4 y M6 & H3, M6 & Azure por canal CSP. Confirmar factibilidad de ExpressRoute con los operadores evaluados. \\
Equipos a bordo, tarjetas y balizas & Curimón & M5--M6 & Piloto y montaje & La unidad de prueba precede al lote. El equipo G26I registra cuatro semanas de fábrica y disponibilidad por confirmar. \\
LTE M2M, Iridium y soporte & audIT & Desde M6; soporte M13 & Antes de su activación & Elegir operador móvil por cobertura. El tiempo de aire se activa por equipo. \\
\end{tabla}

audIT desarrolla el software de negocio y ejecuta la instalación. Compra hardware y licencias de GitLab y SonarQube y contrata nube administrada, conectividad, soporte del fabricante y mantención especializada.

\section{Curva S del proyecto}\label{sec:eco-curva-s}

La curva S cubre los 56 meses de evaluación. Presenta costos acumulados con intereses, facturación acumulada por devengo y saldo de caja financiado. La \figref{eco-curva-s} presenta las series mensuales del modelo consolidado. La amortización del principal se registra en caja. El saldo refleja el momento efectivo de cobro y permite identificar los meses de mayor necesidad de financiamiento.

\figura{economico/costos-venta/curva-s.pdf}{Costos con intereses, facturación y caja financiada acumulados}{eco-curva-s}

\section{Análisis de costos de la solución}\label{sec:eco-costos}

La planilla desglosa personal, costos de terceros, calidad y capacitación, administración y cierre. Los costos de proveedores se registran una sola vez. Las licencias GitLab y SonarQube suman \montoPartida{licencias-d3}. La ejecución auxiliar en Azure pertenece al presupuesto de infraestructura. La compra del hardware de campo y red del mandante se informa por \montoPartida{hardware-mandante} en una hoja separada del flujo de audIT. El pagador de energía y combustible queda sujeto a confirmación.

El presupuesto de riesgos incluye una reserva de infraestructura de \montoPartida{reserva-d4}. Su desembolso se prueba en sensibilidad. Los gastos administrativos incluyen las garantías de fiel cumplimiento, primas, notaría y custodia de fuentes. Las bases asegurables y el overhead incremental siguen en conciliación.

\subsection{Costos frente a venta}\label{sec:eco-venta}

La \tabref{eco-maestro} muestra el costo neto por fase. La comisión de correcto funcionamiento se atribuye a implementación. El modelo estima el desembolso en M21, sujeto a confirmar las condiciones y la fecha de constitución. El costo de las actividades de cierre entre M21 y M24 sigue sujeto a la conciliación de la EDT con las horas y la nómina.

\tablaMontos*{Costo neto por fase}{eco-maestro}{costo-fase-implementacion, costo-fase-operacion, costo-proyecto}

La \tabref{eco-venta} muestra el valor ofertado a partir de esos costos. El margen consolidado sobre ventas netas es 18,50\,\%, dentro del tope de 20\,\%. La métrica de utilidad sobre costo que conserva el Resumen de la plantilla equivale a 22,70\,\%.

\tablaMontos{Valor ofertado por fase}{eco-venta}{implementacion, operacion}

La implementación se factura con 12 hitos de E-25, en proporción 60/40. Los 36 servicios de operación se facturan en M21--M56 y se cobran en M22--M57. La última mensualidad asciende a \montoPartida{cobro-m57} y se concilia como cuenta por cobrar al cierre.

La \tabref{eco-directos-indirectos} resume la asignación vigente de costos y precios. El cuadro maestro de la planilla detalla personal, infraestructura y servicios, licencias, calidad, firma, administración y cierre. La implementación asigna precio en proporción al costo. En operación, el margen se asigna a personal.

\begin{tabla}[ancho=completo]{Asignación de costos y precio por fase}{eco-directos-indirectos}{L{36mm} X X L{17mm}}
Fase y tipo & Costo neto & Venta neta & Margen \\ \midrule
Implementación directa & \montoPartida{cuadro-impl-directo-costo} & \montoPartida{cuadro-impl-directo-venta} & 18,50\,\% \\
Implementación indirecta & \montoPartida{cuadro-impl-indirecto-costo} & \montoPartida{cuadro-impl-indirecto-venta} & 18,50\,\% \\
Operación directa & \montoPartida{cuadro-op-directo-costo} & \montoPartida{cuadro-op-directo-venta} & 18,73\,\% \\
Operación indirecta & \montoPartida{cuadro-op-indirecto-costo} & \montoPartida{cuadro-op-indirecto-venta} & 0,00\,\% \\
\end{tabla}

Los indirectos corresponden a administración, garantías y cierre según su atribución. La conciliación con nómina y asesoría legal debe acreditar cualquier overhead incremental. El IVA referencial por proveedor se concilia con el IVA pagado y el crédito fiscal en la caja financiada.

\section{VAN y TIR de la solución}\label{sec:eco-van}

La evaluación utiliza el flujo libre mensual de la Fila A. El VAN descuenta los flujos de M1 a M56 y considera el flujo de M0. La TIR es la tasa mensual que hace cero el VAN. El cobro de M57 se concilia por separado.

\begin{tabla}{Indicadores de la evaluación financiera}{eco-van}{L{60mm} X}
Indicador & Valor \\ \midrule
Horizonte de evaluación & 56 meses, \form{E-24}{73} \\
Tasa de descuento & 1,0261894\,\% mensual antes de impuestos en el escenario consolidado. Pondera 80\,\% de deuda al 0,9\,\% mensual y 20\,\% de capital propio al 20\,\% efectivo anual supuesto. \\
Valor actual neto (VAN) & \montoPartida{van} \\
Tasa interna de retorno (TIR) & 83,6613\,\% mensual sobre la Fila A de evaluación \\
\end{tabla}

La \tabref{eco-sensibilidad} compara la Fila A mensual a la tasa base. Los shocks reducen los flujos de los primeros meses y la tasa que hace cero su VAN. Por eso se presenta la TIR junto con el VAN y la caja financiada.

\begin{tabla}[ancho=completo]{Sensibilidad del proyecto a precio o calendario fijo}{eco-sensibilidad}{X L{30mm} L{24mm}}
Escenario & VAN, millones de CLP & TIR mensual \\ \midrule
Base & 1.235,0 & 83,66\,\% \\
Costos +5\,\% & 923,7 & 17,32\,\% \\
Cobros −5\,\% & 862,0 & 16,63\,\% \\
H5 cobrado en M13 & 1.230,9 & 83,01\,\% \\
\end{tabla}

El equilibrio nominal es \montoPartida{equilibrio-nominal}. Con costos fijos y escala uniforme de cobros en el calendario base, el equilibrio a VAN cero es \montoPartida{equilibrio-van}. La recuperación descontada sostenida ocurre en M16.

En el caso base, la caja mínima es cero y el uso máximo de la línea es de 45,0 millones de CLP. Si H5 se cobra en M13, la línea alcanza su tope y queda un déficit de \montoPartida{deficit-h5} en M12. Con el desembolso adicional de la reserva de infraestructura, el déficit antes del refuerzo asciende a \montoPartida{deficit-h5-reserva}. El escenario reforzado modela \montoPartida{capital-adicional-h5} en M1, con 80\,\% de crédito y 20\,\% de aporte de socios. La caja mínima calculada es \montoPartida{caja-estres-h5} en M12 frente a una meta de 10 millones de CLP. Si las 72 horas de preparación de CP-SEC-05 se compran íntegramente, la caja simulada baja a \montoPartida{caja-estres-horas}. La viabilidad del escenario depende de la cobertura de horas y de confirmar disponibilidad, giros y comisiones del financiamiento.

\section{Valorización de las cinco innovaciones}\label{sec:eco-innovaciones}

La \tabref{eco-innovaciones} resume el tratamiento de cada innovación. La atribución de personal se contrasta con las horas presupuestadas y su disponibilidad. Los beneficios de Curimón se mantienen separados de los ingresos de audIT.

\begin{tabla}[ancho=completo,fuente=condensada]{Tratamiento de las innovaciones}{eco-innovaciones}{L{32mm} X X}
Innovación & Costo y carga en el modelo & Beneficio para Curimón \\ \midrule
1. Expediente verificable & Firma automatizada presupuestada por \montoPartida{firma-inn1} en M13--M56. Desarrollo y soporte atribuidos a personal. & Monetización pendiente. Medir expedientes, tiempo de preparación y verificaciones. \\
2. Programación de patio & \montoPartida{innovacion-2-inversion} de desarrollo y \montoPartida{innovacion-2-costo} de operación como atribuciones de referencia. Cobertura de horas por conciliar. & Medir visitas y minutos evitados en piloto y compararlos con la línea base. \\
3. Verificación de acople y registro térmico & \montoPartida{innovacion-3-inversion} de integración y \montoPartida{innovacion-3-costo} de operación estimados. Imputación incremental por conciliar. & Medir acoples y eventos térmicos en series independientes. El kilometraje corresponde al criterio de aceptación 26. \\
4. Comodato y recupero de sobreestadías & Los 182 montajes base incluyen 34 comodatos. Las altas adicionales dependen de cupo, capacidad y pagador. & Recupero bruto referencial de \montoPartida{recupero-inn4} en 36 meses a régimen. El neto depende de costos y reparto. \\
5. Prevención de fatiga & \montoPartida{innovacion-5-inversion} de inversión y \montoPartida{innovacion-5-costo} de soporte. Talleres en M4 y calibración en M13--M15. Ejecutores e IVA por confirmar. & Monetización pendiente. Medir consumo y siniestralidad frente a datos históricos. \\
\end{tabla}

Las innovaciones 2 y 3 reconocen cero beneficio monetario en el caso base. El beneficio de la innovación 5 queda pendiente. Los costos atribuidos informados por cada responsable se mantienen identificados hasta completar su cruce con la planilla.

\section{Planilla de cálculo}\label{sec:eco-planilla}

El archivo consolidado reúne el Resumen del mandante, los parámetros, los costos e ingresos, los proveedores, el financiamiento, la curva S y los indicadores. Los cuadros económicos se cotejan con las hojas correspondientes de la misma versión. El flujo mensual presenta el total del período y el saldo acumulado.

\section{Demostración interactiva}\label{sec:eco-demo}

El \href{https://carlosabarza.github.io/PaginaAudIT/#equipo}{sitio de audIT} da acceso al entorno demostrativo. En el sitio, seleccione «Simulador Operativo en Vivo» para iniciar la navegación.
